\section{When adding a site can hurt}
\label{sec:monotonicity}

One might hope that new evidence can only help, so that the design
problem is only about how much a candidate helps. Under the maintained
model this is true.

\begin{theorem}[Monotone improvement under the model]
\label{thm:monotone}
Suppose \cref{ass:model} holds for the pool and for candidate $k$, the
augmented eigenvalue condition holds, and the true coefficient vector
satisfies $\beta\in\restrset$. Then $\feasset'\subseteq\feasset$, both
sets contain $\beta$, and for every target $\ct$,
\begin{equation}
  \width(\ct;\pool\cup\{k\})\;\leq\;\width(\ct;\pool),
  \label{eq:monotone}
\end{equation}
with equality for every target when $\cand{k}\in\affhull(\pool)$. In
particular the augmented program is feasible and its interval still
contains the true target effect $\psi_T$.
\end{theorem}

The assumptions of \cref{thm:monotone} show where it can fail. Every
assumption is a modeling commitment about the new site. The new site
must share the response surfaces $g_a$ and the fixed effects $\beta_a$
of \eqref{eq:response-model}. A site that does not share them, e.g.,
because its implementation or its measurement differs, moves the pooled
normal equations. The normal equations then average conflicting
site-level means, and the identified slice of coefficient space
translates to a compromise that fits no site exactly. The restriction set then cuts the
translated slice at a different cross-section, and the width can move
either way.

\paragraph{A two-site example.}
Take one context coordinate pair, $p=2$, and two source sites with
profiles $c_1=(0,0)$ and $c_2=(1,1)$, the two-site configuration
\citet{boyarsky2026} use to illustrate partial identification. The pool identifies one combination of the
treatment shift, $\delta_1+\delta_2=2$ say, and leaves the direction
$(1,-1)$ free. Impose the box $\shift\in[0,2]\times[-1,1]$ and take a
target that loads on $\delta_1$ alone. The feasible segment is
$\{(\delta_1,2-\delta_1):\delta_1\in[1,2]\}$ and the width is $1$. Now
add a third site on the same line through the profiles whose experiment
implies $\delta_1+\delta_2=0$, in conflict with the pool. Pooling the
two inconsistent moment conditions with equal weight moves the
identified value of $\delta_1+\delta_2$ to the compromise $1$. The new
feasible segment is $\{(\delta_1,1-\delta_1):\delta_1\in[0,2]\}$ and the
width is $2$. A rank-preserving addition doubled the width without changing the
restriction set. Panel (A) of \cref{fig:monotonicity}
draws this example.

\paragraph{What conflict corrupts first.}
Width is not the right measure of the damage. The translated slice can
just as easily cut a narrower cross-section of the restriction set, in
which case the conflicted pool reports a tighter interval around a wrong
location. The precise statement concerns whether the interval covers the true
effect. The translation moves the point the pool identifies, so the
targets that suffer first are the point-identified ones. Their intervals
have width zero, so any movement removes the true effect from the
interval. A wide interval can still cover the true effect after a small
conflict. A point prediction cannot. In the simulation of
\cref{sec:simulation} the conflicted addition changes the worst-case
width by only a few percent in either direction, while the population
interval stops covering the true effect at the span-identified targets
in every replicate. An analyst who sees a surprisingly tight interval after pooling a new
site should treat it as a sign of conflict rather than as an
improvement.

\paragraph{Screens.}
Two checks separate refining additions from corrupting ones.
\begin{itemize}
\item \emph{Rank screen, before the experiment.} Test
  $\cand{k}\notin\affhull(\pool)$. A candidate that passes adds the
  equality constraints \eqref{eq:refinement} and, under the model,
  cannot widen anything (\cref{thm:monotone}). A candidate that fails
  changes no identified quantity, so any movement of the estimated set
  after adding it, beyond sampling noise, is itself evidence of
  conflict.
\item \emph{Compatibility screen, after the experiment.} The candidate's
  data overidentify the in-span coordinates: the projections
  $\Bspan^\top\beta_a$ are estimated by the pool and, once the
  experiment runs, by the candidate site jointly with the pool. Compare
  the pool's minimum-norm fit with the augmented fit on the in-span
  coordinates, in units of their standard errors. A large discrepancy indicates the conflict mechanism above. The check is
an overidentification test of the kind used to detect an invalid
instrument. The forecast
  version of this screen, before the experiment, compares the
  candidate's predicted outcomes under the pool's fit with whatever
  prior information exists about outcomes there.
\end{itemize}
When a screen fails, the analyst should report both analyses, the bounds
with and without the new site, together with the discrepancy between
them. Dropping the new site assumes that the pool is right. Pooling
assumes that the model is right. The screens make the disagreement
visible so the analyst can decide which assumption to keep.
